\section{Technical assumptions and proof overview}
\label{sec:proof-overview}

\subsection{Technical Assumptions}
\label{subsec:tracking-assumptions}

We state the conditions used in the main theorem and proof sketch;
\Cref{subsec:general-argument-bound-settings} gives their full analytic formulation,
and \Cref{tab:assumption-evidence-map} records how each model section relates to
the theorem's assumptions and conclusions.
The conditions below are verified for
\Cref{ex:variance-cosine-curvature-recovery} in
\Cref{lem:canonical-comparison-assumptions}.
\begin{assumption}[Structured moment representation]
\label{assumption-main-general-moment-smooth-compressed}
The full conditions of \Cref{assumption-structured-moment-smoothness} hold.
We use \(F(\rho)=G(M(\rho))\), with
\(M(\rho):=(\int\phi_m\,d\rho)_{m\ge1}\), as defined in
\Cref{def:moment-objective}.
The features have controlled polynomial growth through third spatial
derivatives, dominated by the decay of the moment weights
\(a_m=e^{-c_am^{\kappa_{\rm dec}}}\), with \(c_a>0\) and
\(\kappa_{\rm dec}>1\); see
\Cref{eq:weighted-moment-space,assumption-structured-moment-smoothness}.
The functional \(G\) is \(C^3\) on one open weighted-moment domain
containing all finite-scale feasible moment regions, with a uniform
weighted third-derivative bound.
All structural parameters, including the anchor derivative norms and this
uniform bound, are fixed independently of \(N,T,p\).
\end{assumption}

The localized drift and Hessian bounds derived in
\Cref{lem:conditional-local-coefficient-growth,lemma-localized-moment-smoothness-consequences}
may depend on the spatial localization radius \(D\).

\begin{assumption}[Regularity of the coupled particle systems]
\label{assumption-main-coupled-particle-regularity}
Fix \(T>0\). The coupled particle systems
\(\mathbf X_t=(X_t^1,\dots,X_t^N)\) and
\(\mathbf Y_t=(Y_t^1,\dots,Y_t^N)\) satisfy the following conditions uniformly
on \([0,T]\). The \(O(1)\) bounds below are uniform in \(N\) over the
horizons considered in \Cref{theorem-main-r_t}.

\begingroup\raggedright
\noindent
\textbf{(i) Unconditional moment bound.}
For fixed $\alpha\in(0,2]$,
\(\sup_{t\le T,q\ge1}q^{-1/\alpha}
\max\{(\mathbb E\|X_t^1\|^q)^{1/q},
       (\mathbb E\|Y_t^1\|^q)^{1/q}\}=O(1)\).\par
\endgroup

\noindent
\textbf{(ii) Temporal regularity.}
All particle paths are right-continuous with left limits at the jumps and
share a deterministic \(O(1)\)-Lipschitz bound on every ODE branch, almost surely.

\noindent
\textbf{(iii) Non-spikiness of $r_t$.}
The coupling error \(r_t\) from \Cref{subsec:tracking-metric} is assumed,
almost surely, not to concentrate on a small number of particles: for some
\(A_{\mathrm{sp}} = O(1)\),
\(\frac1N\sum_{i=1}^N\|X_t^i-Y_t^i\|^3 \le A_{\mathrm{sp}} r_t^3\;
    (0\le t\le T)\).
This condition controls the cubic Taylor remainder in the particle positions by \(O(r_t^3)\);
see \Cref{lemma-directional-point-slot-error,eq-directional-point-slot-bound}.
In \Cref{ex:variance-cosine-curvature-recovery}, all paired errors coincide, giving
\(A_{\mathrm{sp}}=1\) (\Cref{lem:canonical-comparison-assumptions}(v)).
\end{assumption}

\begin{assumption}[Prescribed schedule and population-first Gaussian jumps]
\label{assumption-main-jump-time}
The common times $(t_\ell)_{\ell=1}^{\NofJ}$ and amplitudes $\eta_\ell\ge0$
are deterministic, with $\NofJ=O(1)$ in \Cref{theorem-main-r_t}.
The initial law \(\mu_0\) is deterministic; the common initial positions
\(X_0^i=Y_0^i\), \(i=1,\ldots,N\), are sampled independently from \(\mu_0\).
Population innovations are independent of that sample and fresh at each
jump. The population Gaussian fields and the empirical transitions satisfy
Assumption~\ref{assumption-common-seed}.
\end{assumption}

\subsection{Proof sketch}
\label{subsec:proof-sketch-after-assumptions}

We explain how the ODE and jump comparisons produce the response factors
in \Cref{def:curvature-responses}. The calculations below are schematic:
\(\simeq\) retains the leading linearized drift, and \(\lesssim\) suppresses
localized coefficients and Gaussian probability factors quantified in the appendix.

The reference particles are i.i.d. with law \(\mu_t\) conditional on
\(\mathcal G_{\rm pop}\) by
\Cref{lemma-gp-jump-implies-conditional-iid,lem:population-first-finite-coupling},
enabling concentration (background: \citet{sznitman1991,fournier2015rate}).
We use \(\Gamma_t^N\) and \(\widetilde\Gamma_t^N\) from
\Cref{subsec:general-argument-overview} for drift/Hessian fluctuations
and covariance-sampling error at jumps, respectively.

\maintextheading{ODE intervals (\Cref{subsec:general-argument-bound-r_t-ode-continuous}).}
Write \(\Delta_t:=\mathbf X_t-\mathbf Y_t\), with particle blocks indexed by \(i\),
so that \(r_t=N^{-1/2}\|\Delta_t\|_{\mathbb R^{dN}}\).
For \(r_t>0\), the chain rule and the ODEs in
\cref{eq:particle-flow,eq:reference-flow} give the equalities below. After adding and subtracting the reference empirical drift,
then linearizing and comparing Hessians
(\Cref{lemma-mean-value-theorem,lemma-bound-r_t-given-eta-eps}), we informally have
\[
 r_t\dot r_t=\frac12\frac{\mathrm d}{\mathrm dt}r_t^2
 =-\frac1N\sum_{i=1}^N\bigl\langle\Delta_t^i,
 \underbrace{\nabla_WF(\rho_{\mathbf X_t},X_t^i)-\nabla_WF(\mu_t,Y_t^i)}_{
 \simeq\,(\HessianMatrices_t^Y\Delta_t)^i\;\text{(Taylor)}}\bigr\rangle
 \simeq-\frac1N\Delta_t^\top\HessianMatrices_t^Y\Delta_t.
\]
By \Cref{def:curvature-responses}, the quadratic form is bounded above by
\(-\kappa_t^{Y,N}r_t^2\); dividing by \(r_t\) gives the curvature term
\(-\kappa_t^{Y,N}r_t\).
Restoring the errors, on the localized good events and while \(r_t\le R\),
\Cref{lemma-common-empirical-moment-residual-directional-rt},
\eqref{eq-rt-directional-stopped}, gives
\(\dot r_t\le-\kappa_t^{Y,N}r_t+c_DRr_t+\Gamma_t^N(1+r_t)\).
Here \(c_D\) is defined in \eqref{eq-Csp-definition}; the same lemma
also handles \(r_t=0\).

\maintextheading{Jump times (\Cref{subsec:general-argument-r_t-jump-errors}).}
Using the empirical and aligned population factors \(M_X,B_{\rm pop}\)
and shared Gaussian vector \(\mathbf g_\ell^N\) from
\Cref{subsec:auxiliary-particles}, we substitute the jump updates in
\cref{eq:particle-jump,eq:reference-jump} and apply the triangle inequality:
\[
 r_{t_\ell}
 =\frac1{\sqrt N}\bigl\|\Delta_{t_\ell-}
 +\eta_\ell(M_X-B_{\rm pop})\mathbf g_\ell^N\bigr\|_{\mathbb R^{dN}}
 \le r_{t_\ell-}
 +\frac{\eta_\ell}{\sqrt N}
 \bigl\|(M_X-B_{\rm pop})\mathbf g_\ell^N\bigr\|_{\mathbb R^{dN}}.
\]
Let \(M_Y\) be the empirical Hessian factor at the reference configuration
just before the jump. We split
\(M_X-B_{\rm pop}=(M_X-M_Y)+(M_Y-B_{\rm pop})\).
Smoothness controls the contribution of the first difference by a multiple
of \(r_{t_\ell-}\). For the remaining mismatch, Gaussian concentration and
the aligned coupling give schematically, with all time-\(t\) quantities
evaluated at \(t_\ell-\),
\[
 \frac{\|(M_Y-B_{\rm pop})\mathbf g_\ell^N\|_{\mathbb R^{dN}}}{\sqrt N}
 \lesssim
 \underbrace{
 \sqrt{\frac1N
 \|M_Y^2-B_{\rm pop}B_{\rm pop}^{\top}\|_{*;\,\mathbb R^{dN}\to\mathbb R^{dN}}}
 }_{\text{square root of covariance mismatch}}
 \lesssim \sqrt{\Gamma_t^N}+\sqrt{\widetilde\Gamma_t^N}.
\]
Here \(*\) denotes the nuclear norm. The \((i,j)\)-block of the covariance
mismatch is \(K_{\rho_{\mathbf Y_t}}(Y_t^i,Y_t^j)-K_{\mu_t}(Y_t^i,Y_t^j)\).
Using the kernel definition in \Cref{subsec:population-pwgf}, the difference
\(K_{\rho_{\mathbf Y_t}}(x,y)-K_{\mu_t}(x,y)\) consists schematically of
\[
 \underbrace{\text{Hessian approximation}}_{\text{as in the ODE step}}
 +\int\nabla_W^2F(\mu_t;x,z)\,
 \nabla_W^2F(\mu_t;z,y)\,
 (\rho_{\mathbf Y_t}-\mu_t)(dz).
\]
Expanding the difference of the two Hessian products bounds the first term
linearly in the Hessian approximation error, with localized coefficients suppressed.
The second term is an empirical integration error.
These comparisons are established in
\Cref{lemma-aligned-finite-square-root,lemma-kernel-mismatch-bound,lem:jump-bound-common-seed}.
Combining the two contributions yields
\(r_{t_\ell}\lesssim(1+\eta_\ell)r_{t_\ell-}
+\eta_\ell(\sqrt{\Gamma_{t_\ell-}^N}+\sqrt{\widetilde\Gamma_{t_\ell-}^N})\):
a jump amplifies past error and adds a new error.
The square-root dependence explains why an \(N^{-1/2}\) covariance-sampling
bound contributes \(N^{-1/4}\) to the jump estimate.
For background on Gaussian covariance alignment, see \citet{masarotto2019procrustes}.

\maintextheading{From local errors to response factors.}
Between jumps, the homogeneous equation
\(\dot r_t=-\kappa_t^{Y,N}r_t\) has solution
\(r_t=G_\kappa^Y(t,s)r_s\), where
\(G_\kappa^Y(t,s)=\exp(-\int_s^t\kappa_u^{Y,N}\,du)\).
Hence an error created at time \(s\) receives this weight at time \(t\).
Starting from \(r_0=0\), integrating the continuous forcing and summing
the jump forcing gives an integral and a sum, respectively.
Taking uniform residual bounds outside them yields schematically
\[
 \textstyle r_t\lesssim
 (\sup_{s\le T}\Gamma_s^N)
 \underbrace{\textstyle\int_0^tG_\kappa^Y(t,s)\,ds}_{\mathfrak C_t^Y}
 +(\sqrt{\sup_{s\le T}\Gamma_s^N}
       +\sqrt{\max_{\ell\le\NofJ}\widetilde\Gamma_{t_\ell-}^N})
 \underbrace{\textstyle\sum_{\ell:t_\ell\le t}\eta_\ell
                  G_\kappa^Y(t,t_\ell)}_{\mathfrak J_t^Y}.
\]
The maximum is zero if there are no jumps.
The recursive comparison in \Cref{lem:piecewise-gronwall-forced-jumps},
\cref{eq:piecewise-recursive-bound,eq:tracking-stopped-propagator-bound},
retains the extra nonlinear amplification and the products of jump multipliers.

\maintextheading{Residual bounds and closure.}
Conditioning only on \(\mathcal G_{\rm pop}\), moment localization and
weighted-feature covariance concentration give, with probability at least
\(1-p\) and \(\Xi=\Xi_{N,T,p}\) from \Cref{theorem-main-r_t},
\[
 \sup_{t\le T}\Gamma_t^N
 \le C_\varepsilon\Xi^\varepsilon(N^{-1/2}+T/N),\qquad
 \max_{\ell\le\NofJ}\widetilde\Gamma_{t_\ell-}^N
 \le C_\varepsilon\Xi^\varepsilon N^{-1/2};
\]
see \Cref{lemma-uniform-sup-gamma-bound-convenient,cor:tilde-gamma-jump-time-union-bound}.
Under the theorem's conditions,
\Cref{thm:fixed-schedule-high-prob-r-v15,cor:fixed-schedule-high-prob-r-rate-v15}
close the localized estimate and absorb the extra amplification into
\(C_\varepsilon\Xi^\varepsilon\).
\Cref{lem:functional-gap-fixed-time-v15,eq:functional-gap-fixed-time-v15}
give the objective bound.
